\documentclass[12pt]{article}
\usepackage[utf8]{inputenc}
\usepackage[T1]{fontenc}
\usepackage{hyperref}
\title{Meskwaki Language Resources for K-12 Curriculum Development}
\author{Compiled by Claude Sonnet 5.5 (claude-sonnet-5-5) for Andr\'as Kornai}
\date{}
\begin{document}
\maketitle

\section{Introduction}
Meskwaki (ISO 639-3 sac) is an endangered Algic language spoken at the Meskwaki Settlement in Iowa, USA. The community runs its own language preservation work and a school with a language programme.

\section{Primary Educational Resources}
\begin{itemize}
\item Meskwaki Settlement School (Meskwakiatoweni page). Search snippets say the school is working toward full Meskwaki immersion. They also say the early childhood and pre-K programmes became immersion in 2014. These pages may be dated, so confirm the current status with the school: \url{https://msswarriors.org/school/meskwakiatoweni/}
\end{itemize}

\section{Linguistic and Reference Materials}
\begin{itemize}
\item Smithsonian note on the Meskwaki dictionary by Ives Goddard and Lucy Thomason. It runs to nearly 500 pages and draws on archival writings by native Meskwaki speakers: \url{https://naturalhistory.si.edu/research/anthropology/programs/recovering-voices/outreach/news-research-updates/meskwaki-dictionary-mother-languages}
\end{itemize}

\section{Religious and Community Resources}
\begin{itemize}
\item Meskwaki Nation Language Preservation Department. Its mission includes developing classroom resources and a repository of language materials: \url{https://www.meskwaki.org/language-preservation/}
\item Our Iowa Heritage page on the Meskwaki language: \url{https://ouriowaheritage.com/our-iowa-heritage-meskwaki-language/}
\end{itemize}

\section{Digital and Audio Resources}
The Language Preservation page describes (according to the search summary) a Meskwaki language app for iOS (2015) and Android (2016), eleven iBooks, card games and language booklets. Adult classes also run, including over Zoom. I did not verify the app store listings.

\section{Curriculum Development Context}
Meskwaki has a community-run language department, a school with immersion ambitions, and a major scholarly dictionary. The next step is to contact the Language Preservation Department and the school directly about their materials.
\end{document}